\documentclass[../../../include/open-logic-section]{subfiles}

\begin{document}

\olfileid{sth}{replacement}{extrinsic}
\olsection[外在考量]{关于替换公理的外在考量}

我们首先考虑一种为替换公理提供\emph{外在}证成的尝试。Boolos（布尔）提出了如下一种。


\begin{quote}
  [\ldots] 采用替换公理的理由非常简单：它们有许多可取的结果，并且（显然）没有不良后果。除了关于迭代概念的定理之外，这些结果还包括一个或许不够理想但令人满意的无穷数理论，以及一个非常可取的结果，该结果证成了良基关系上的归纳定义。
  \citep[229]{Boolos1971}
\end{quote}

Boolos 观点的要旨在于，我们应当通过其成果来证成替换公理。

而他提及的具体成果，正是我们在前几章中讨论过的内容。

替换公理使我们能够证明 von Neumann（冯·诺伊曼）序数是良序型概念的出色代表（这就是我们“或许不够理想但令人满意的无穷数理论”）。

替换公理还允许我们定义 $V_\alpha$，确立秩的概念，并证明 $\in$-归纳（这相当于我们的“关于迭代概念的定理”）。

最后，替换公理允许我们证明超限递归定理（这就是“良基关系上的归纳定义”）。

这些确实是可取的结果。

但是，这些可取的结果足以\emph{证成}替换公理吗？\emph{不能}。

或者至少，不能直接证成。

这里有一个简单的问题。

尽管我们列出了替换公理的一些可取结果，但其中许多结果本可以通过其他途径获得。

不过，这一点并未得到应有的普及，所以我们有必要稍作停顿来解释一下情况。

存在一种简单的集合论，称为层理论（Level Theory），简称 $\LT$。\footnote{$\LT$ 的最初版本由 \citet{Montague1965} 和 \citet{Scott1974} 提出；\citet{Potter2004} 对此进行了简化，并给出了专著篇幅的论述；而 \citet{ButtonLT1} 最近进一步简化了 $\LT$。} $\LT$ 的公理只有外延公理、分离公理，以及每个集合都是某个\emph{层}的子集这一断言，其中“层”被巧妙地定义，使得这些层的表现如同我们的老朋友 $V_\alpha$ 一样。

因此 $\ZF$ 能证明 $\LT$；但 $\LT$ 比 $\ZF$ \emph{弱得多}。

事实上，$\LT$ 无法给出对公理、幂集公理、无穷公理或替换公理。

令 $\Zr$ 为在 $\LT$ 中加入无穷公理和幂集公理后得到的系统；这也蕴含了对公理，因此，$\Zr$ 至少与 $\Z$ 一样强。

但是，事实上 $\Zr$ 严格强于 $\Z$，因为它添加了每个集合都有秩这一断言（因此我建议将其称为 $\Zr$）。

确实，$\Zr$ 能够给出：一个完全令人满意的序数理论；
将层级结构划分为良序阶段的结果；$\in$-归纳的证明；以及超限递归定理的一个\emph{版本}。

简而言之：尽管 Boolos 不知道这一点，但他所提到的所有可取结果，本都无需替换公理即可达到；他只需使用 $\Zr$ 而非 $\Z$。

（考虑到这一切，为什么我们还要遵循传统路径，向你讲授 $\ZF$，而不是 $\LT$ 和 $\Zr$ 呢？

有两个原因。

第一：纯粹由于历史原因，从 $\LT$ 开始是相当非标准的；我们希望让你具备阅读更标准的集合论文献的能力。

第二：当你准备好理解 $\LT$ 和 $\Zr$ 时，你可以直接阅读 \citealt{Potter2004} 和 \citealt{ButtonLT1}。）

当然，由于 $\Zr$ 严格弱于 $\ZF$，有些结果是 $\ZF$ 能证明而 $\Zr$ 无法判定的。

因此，人们可以试图通过指出其中某个结果，来为替换公理提供外在证成。

但是，一旦你知道如何使用 $\Zr$，就很难找到许多同时满足以下条件的例子：（a）由替换公理判定但无法通过其他方式判定，并且（b）直觉上为真。（有关此点的更多内容，参见 \citealt[\S13.2]{Potter2004}。）

归根结底：

要为替换公理提供一个有说服力的外在证成，需要找到一个没有替换公理就\emph{无法}实现的结果。

而这并非易事。

让我们考虑一个进一步的问题，它出现在任何试图为替换公理提供纯粹外在证成的尝试中。（这个问题或许比第一个更根本。）Boolos 并不仅仅指出替换公理有许多可取的结果。

他还指出替换公理“（显然）没有”不良后果。

但是，这个括号中的限定词“显然”，无疑是绝对至关重要的。

回想一下我们是如何走到这一步的：朴素概括公理陷入了不一致，而我们通过接受集合的累积迭代概念来应对这一不一致。

这一概念自带了一套叙事，我们希望它能向我们保证其一致性。

但是，如果我们无法从那套叙事内部证成替换公理，那么我们（迄今）就没有理由相信 $\ZF$ 是一致的。

或者说：除了至今尚无人发现矛盾这一（或许只是偶然的）事实之外，我们没有理由相信 $\ZF$ 是一致的。

正是在这个意义上，Boolos 的评论似乎可以归结为：“（显然）$\ZF$ 是一致的”。

我们应当要求比这更强的一致性保证。

这个问题将影响任何对替换公理的\emph{纯粹}外在证成的尝试，即任何仅仅以 $\ZF$ 的（已知）结果来表述的证成。

因此，我们将希望寻找替换公理的一种\emph{内在}证成，即一种能表明我们关于集合的叙事在某种意义上“早已”蕴含了替换公理的证成。

\end{document}